\chapter{Family--defect flat-holonomy realization gate}

Projector-supercurvature gate оставил один algebraic post-processing step:
после выбора \(P_a\) permutation \(C_{a,\nu}\) назначалась отдельно. Здесь
строится явная connection one-form, чья holonomy точно равна этой
permutation и использует тот же projector field.

\section{Канонический orientation generator}

Пусть
\begin{equation}
 u=\frac12(1,1,1,1)^T
\end{equation}
--- нормированный affine singlet, а \(h_a=\operatorname{diag}H_a\) ---
нормированная triplet axis, выбранная equation \(Q(H_a)=0\). Введём
\begin{equation}
 \Omega(h_a)_{bc}
 =\sum_{d,e=1}^4\epsilon_{bcde}u_d(h_a)_e.
 \label{eq:orientation-generator}
\end{equation}
Этот operator задаётся только canonical singlet, orientation tensor и
curvature-selected axis. Он удовлетворяет
\begin{align}
 \Omega^T&=-\Omega,\\
 \Omega u&=0,\\
 \Omega h_a&=0,\\
 \Omega^2&=-\left(I_4-uu^T-h_ah_a^T\right).
\end{align}
Следовательно, \(\Omega\) является unit complex structure на двумерной
плоскости, ортогональной singlet и выбранной axis.

Так как \(H_a\) имеет один distinguished eigenvalue и тройной eigenvalue
на complement, получаем также
\begin{equation}
 [\Omega(h_a),H_a]=0.
 \label{eq:connection-field-commutator}
\end{equation}

\section{Flat connection и точная holonomy}

На boundary circle \(\gamma\) длины \(L\) определим
\begin{equation}
 \mathcal A_{a,\nu}
 =-\frac{2\pi\nu}{3L}\Omega(h_a)\,ds,
 \qquad \nu=\pm1.
 \label{eq:projector-flat-connection}
\end{equation}
Connection постоянна и одномерна, поэтому flat. Её holonomy равна
\begin{equation}
 \operatorname{Hol}_\gamma(\mathcal A_{a,\nu})
 =\exp\!\left[-\nu\frac{2\pi}{3}\Omega(h_a)\right].
\end{equation}

Exponential фиксирует \(u\) и \(h_a\), а на transverse plane поворачивает
на угол \(-2\pi\nu/3\). Exhaustive calculation даёт exact identity
\begin{equation}
 \operatorname{Hol}_\gamma(\mathcal A_{a,\nu})=C_{a,\nu}
 \label{eq:holonomy-equals-three-cycle}
\end{equation}
для всех четырёх projectors и обоих winding sectors. Максимальный
numerical residual меньше \(1.8\times10^{-15}\).

\begin{theorem}[Flat-holonomy realization]
Каждый из восьми \(S_4\) three-cycles является точной holonomy
connection~\eqref{eq:projector-flat-connection}, построенной из того же
rank-one projector, который решает finite curvature equation
\(Q(H)=0\). Следовательно, cycle больше не является внешним
post-processing assignment.
\end{theorem}

\section{Twisted covariance}

Для \(g\in S_4\) orientation generator преобразуется как
\begin{equation}
 g\Omega(h_a)g^{-1}
 =\operatorname{sgn}(g)\Omega(h_{g(a)}).
\end{equation}
Если winding является orientation-line variable,
\begin{equation}
 \nu\longmapsto\operatorname{sgn}(g)\nu,
\end{equation}
то connection и holonomy ковариантны:
\begin{align}
 g\mathcal A_{a,\nu}g^{-1}
 &=\mathcal A_{g(a),\operatorname{sgn}(g)\nu},\\
 gC_{a,\nu}g^{-1}
 &=C_{g(a),\operatorname{sgn}(g)\nu}.
\end{align}
Все 192 cases проходят с нулевым connection-covariance residual.

\section{Один constrained superconnection saddle}

Объединим ordinary connection и finite odd field символически как
\begin{equation}
 \mathbb A_{a,\nu}=d+\mathcal A_{a,\nu}+T(H_a),
\end{equation}
где shifted square \(T(H)^2-\frac13I\) воспроизводит
\(Q(H)\). На selected branch одновременно выполнены
\begin{equation}
 F_{\mathcal A}=0,
 \qquad
 D_{\mathcal A}H_a=0,
 \qquad
 Q(H_a)=0.
\end{equation}
Второе равенство следует из
equation~\eqref{eq:connection-field-commutator}. Поэтому finite curvature
и boundary holonomy совместимы как блоки одного constrained
superconnection saddle.

После ограничения на standard triplet integrated connection имеет rank
два и одномерное kernel. Тем самым тот же operator даёт и exact-one
Majorana selector.

Определение holonomy superconnection через parallel transport является
стандартной конструкцией; см. \path{arXiv:0912.0249}. Flat defect
connections с holonomies в фиксированных conjugacy classes также имеют
прямые аналоги, например \path{arXiv:1304.1372}. Эти источники
подтверждают архитектуру, тогда как explicit formula
\eqref{eq:orientation-generator} является результатом project affine
carrier.

\section{Constitutive-action criterion}

Положительный theorem ещё использует constitutive saddle relation
\begin{equation}
 \mathcal A
 =-\frac{2\pi\nu}{3L}\Omega(h)\,ds.
 \label{eq:constitutive-holonomy-relation}
\end{equation}
В пределах этой главы relation канонична и ковариантна, но ещё не получена
как Euler--Lagrange equation одного local, topological или spectral action.
Поэтому здесь её зависимость от \(H\) всё ещё имеет статус ansatz.

\begin{nogocriterion}[Constitutive-action gate]
Полное boundary-superconnection замыкание допускается только если
relation~\eqref{eq:constitutive-holonomy-relation} следует из вариации
одного заранее фиксированного action с тем же trace normalization, что и
\(\Tr Q(H)^2\). Независимый multiplier, введённый только для навязывания
этой relation, не считается выводом.
\end{nogocriterion}

Следующий cubic-root variational gate проходит этот критерий в фиксированном
nonzero unit-winding sector. Для frame
\(W=\exp[(\varphi/3)\Omega(H)]\) variation единого positive action по
connection даёт
\begin{equation}
 \mathcal A_s=-\partial_sW\,W^T
 =-\frac{\partial_s\varphi}{3}\Omega(H),
\end{equation}
а integrated holonomy равна \(C_{a,\nu}\) независимо от phase profile.

Таким образом, algebraic holonomy gap и constitutive-action gap закрыты.
Открытым остаётся более глубокий dynamical вопрос: вывести nonzero
cubic-root condensate и его bundle embedding из полного finite graded
parent action.

Машинный ledger сохранён в
\path{s2t_v4_family_defect_holonomy_realization_gate_results.json};
генератор ---
\path{s2t_v4_family_defect_holonomy_realization_gate.py}.